%% 附录

\begin{thesisappendix}
\section{数据集及构建补充材料}
\label{app:dataset-construction}

本附录补充 UAV-DualCog 测试基准的数据规模、任务结构、构建流程、评测协议与能力优化实现细节。相关资源已公开发布，便于对测试基准、代码实现和数据样本进行复核。

\subsection{公开发布与复现资源}
\label{app:release-resources}

UAV-DualCog 的项目站点、代码仓库和数据集已分别发布，三类资源的作用见\autoref{tab:app-release-resources}。其中，项目站点用于集中展示测试基准、评测协议、分析结果和使用说明；代码仓库提供数据构建、评测和能力优化的主要实现；ModelScope 数据集页面提供测试基准的数据下载入口。

\begin{table}[htbp]
    \centering
    \caption{UAV-DualCog 公开发布资源}
    \label{tab:app-release-resources}
    \small
    \begin{tabularx}{\textwidth}{@{}U{2.5cm}X X@{}}
        \toprule
        资源类型 & 地址 & 作用说明 \\
        \midrule
        测试基准站点 & \url{https://uav-dualcog.lozumi.com/} & 展示基准设计、评测协议、排行榜、实验分析和使用说明。 \\
        代码仓库 & \url{https://github.com/SmartDianLab/UAV-DualCog} & 提供数据构建工具链、评测脚本、训练实现和复现实验所需配置。 \\
        数据集页面 & \url{https://www.modelscope.cn/datasets/Lozumi/UAV-DualCog} & 发布 UAV-DualCog 数据集资源，用于下载测试基准及相关数据文件。 \\
        \bottomrule
    \end{tabularx}
\end{table}

\begin{table}[htbp]
    \centering
    \caption{UAV-DualCog 测试基准发布规模补充统计}
    \label{tab:app-dataset-scale}
    \small
    \begin{tabularx}{\textwidth}{@{}U{3.0cm}U{3.0cm}X@{}}
        \toprule
        统计项 & 数值 & 说明 \\
        \midrule
        测试场景数 & 12 & 来自与训练场景隔离的 AerialVLN 仿真环境 \\
        有效地标数 & 512 & 经过多视角复核和语义标注后的发布地标 \\
        图像任务样本数 & 4096 & 覆盖四类图像任务与两条双认知能力主线 \\
        视频任务样本数 & 2048 & 覆盖飞行行为识别与地标可见性推理 \\
        图像发布分辨率 & DCI 4K / JPEG & 保留高分辨率地标细节，评测时可按模型能力压缩输入 \\
        视频发布规格 & 1080P / 10 FPS / MP4 & 支持行为识别、可见性计数和时间区间定位 \\
        \bottomrule
    \end{tabularx}
\end{table}

\begin{table}[htbp]
    \centering
    \caption{UAV-DualCog 六类任务与证据形式}
    \label{tab:app-task-matrix}
    \footnotesize
    \setlength{\tabcolsep}{3pt}
    \begin{tabularx}{\textwidth}{@{}U{2.8cm}U{1.4cm}U{2.0cm}U{4.5cm}U{2.4cm}@{}}
        \toprule
        任务类型 & 媒介 & 认知分支 & 输入形式 & 输出与证据 \\
        \midrule
        地标参照位置推理 & 图像 & 自我状态认知 & 地标参考图 + 当前观测图 & 选项 + 目标框 \\
        未来观察预测 & 图像 & 自我状态认知 & 地标参考图 + 候选未来视图 & 选项 + 目标框 \\
        自我参照位置推理 & 图像 & 环境情况认知 & 当前观测图 + 地标描述 & 选项 + 目标框 \\
        地标驱动动作决策 & 图像 & 环境情况认知 & 当前观测图 + 地标描述 & 选项 + 目标框 \\
        飞行行为识别与时间定位 & 视频 & 自我状态认知 & 第一视角飞行视频 & 行为选项 + 时间区间 \\
        地标可见次数与时间定位 & 视频 & 环境情况认知 & 参考图 + 第一视角视频 & 可见次数 + 时间区间 \\
        \bottomrule
    \end{tabularx}
\end{table}

\begin{table}[htbp]
    \centering
    \caption{测试基准构建工具链主要入口}
    \label{tab:app-toolchain}
    \scriptsize
    \begin{tabularx}{\textwidth}{@{}U{1.7cm}U{3.8cm}X@{}}
        \toprule
        阶段 & 主要入口 & 作用说明 \\
        \midrule
        阶段1 & \makecell[c]{\texttt{stage1\_collect\_}\\\texttt{pcd.py}} & 调用 AerialVLN 模拟器进行场景扫描，采集语义分割、位姿和点云信息，并融合为场景级几何基础。 \\
        阶段2 & \makecell[c]{\texttt{stage2\_landmark\_}\\\texttt{label.py}} & 基于实例点云生成地标候选，支持人工复核、语义标注和有效地标清单导出。 \\
        阶段3 & \makecell[c]{\texttt{stage3\_generate\_}\\\texttt{traj.py}\\\texttt{stage3\_task\_}\\\texttt{suite.py}} & 生成飞行轨迹、视频资产、复合行为和原子行为时间监督，并导出视频任务。 \\
        阶段4 & \makecell[c]{\texttt{stage4\_qa\_generate\_}\\\texttt{and\_eval.py}} & 围绕地标中心参考视图和无人机观测图生成图像问答任务，并导出目标框监督。 \\
        全流程 & \makecell[c]{\texttt{task\_pipeline.py}\\\texttt{pipeline\_common.py}} & 组织阶段间输入输出、路径管理、任务配置和批处理运行。 \\
        \bottomrule
    \end{tabularx}
\end{table}

\subsection{构建工具链可视化工作台界面}

为补充正文第3章对阶段2、阶段3和阶段4实现过程的说明，本节收录数据构建 Web 工作台中的代表性界面截图。该工作台面向人工复核、任务定义、样本浏览、实验配置和指标检查等交互环节；大规模样本构建与导出仍由命令行工具链统一调度。下列截图覆盖不同构建环节的主要职责，用于说明可视化复核与质量控制方式，而不作为实验结果数据来源。

\begin{figure}[!htbp]
    \centering
    \includegraphics[width=\textwidth]{fig_app_stage2_review_workspace.png}
    \caption{阶段2地标复核与语义标注工作台界面}
    \label{fig:app-stage2-review-workspace}
\end{figure}

\autoref{fig:app-stage2-review-workspace} 展示了阶段2地标复核界面。左侧为候选实例列表和保留状态，中间区域同时展示点云主视图、点云俯视图、RGB 俯视图以及多方向 RGB 侧视图，右侧为实例属性、自动标注结果和人工修订字段。该界面使标注者能够在同一页面中检查实例类别、点云形态、遮挡情况、主视图选择、方向标签和语义描述，从而支撑有效地标清单的质量控制。

\begin{figure}[!htbp]
    \centering
    \includegraphics[width=\textwidth]{fig_app_stage3_behavior_library.png}
    \caption{阶段3飞行行为库配置工作台界面}
    \label{fig:app-stage3-behavior-library}
\end{figure}

\autoref{fig:app-stage3-behavior-library} 展示了阶段3的飞行行为库界面。左侧列出复合行为和原子行为定义，右侧展示所选行为的参数范围、默认值、相机模式和生成说明。该界面使复合行为、原子行为和轨迹参数之间的关系显式化，有助于检查层次化飞行模式是否符合视频任务设计，并形成可复用的行为模板。

\begin{figure}[!htbp]
    \centering
    \includegraphics[width=\textwidth]{fig_app_stage3_mission_generation.png}
    \caption{阶段3视频任务 Mission Generation 工作台界面}
    \label{fig:app-stage3-mission-generation}
\end{figure}

\autoref{fig:app-stage3-mission-generation} 展示了阶段3视频任务构建的任务生成界面。该界面围绕场景、地标选择、复合行为模板、原子行为序列、轨迹参数和媒体导出设置组织任务构建信息，用于把阶段2得到的有效地标资产转换为可执行飞行任务。该界面同时记录任务选择、轨迹配置、预览复核和视频导出相关信息，并与后续清单核查和指标统计环节衔接。

\begin{figure}[!htbp]
    \centering
    \includegraphics[width=\textwidth]{fig_app_stage3_dataset_browser.png}
    \caption{阶段3视频任务样本浏览与证据核查界面}
    \label{fig:app-stage3-dataset-browser}
\end{figure}

\autoref{fig:app-stage3-dataset-browser} 展示了阶段3的视频任务样本浏览界面。左侧按照样本编号、任务类型、难度和地标编号列出任务清单，右侧给出样本级元数据、地标参考图像、飞行描述、视频规格和证据字段。该界面用于检查视频任务是否同时包含语义问题、地标视觉证据、行为描述和时间监督信息，避免任务清单与媒体文件、参考图像或标注区间脱节。

\begin{figure}[!htbp]
    \centering
    \includegraphics[width=\textwidth]{fig_app_stage3_metrics.png}
    \caption{阶段3视频任务指标汇总工作台界面}
    \label{fig:app-stage3-metrics}
\end{figure}

\autoref{fig:app-stage3-metrics} 展示了阶段3的指标汇总界面。该界面按照任务类型、飞行模式、难度和模型运行结果聚合主要指标，并保留实验矩阵用于交叉检查。它对应正文中视频任务评测所需的复合行为识别、原子行为识别、地标可见次数和时间区间定位等指标，是从样本级输出过渡到整体实验分析的重要核查入口。

\begin{figure}[!htbp]
    \centering
    \includegraphics[width=\textwidth]{fig_app_stage4_task_generation.png}
    \caption{阶段4图像 QA 任务生成工作台界面}
    \label{fig:app-stage4-task-generation}
\end{figure}

\autoref{fig:app-stage4-task-generation} 展示了阶段4图像任务生成界面。该界面按场景、任务类型、难度、地标类别和采样数量组织图像 QA 样本构建信息，右侧同时给出场景可用地标、任务可用性和实验摘要。该设计便于在样本构建前检查地标是否支持四类图像任务。

\begin{figure}[!htbp]
    \centering
    \includegraphics[width=\textwidth]{fig_app_stage4_data_browser.png}
    \caption{阶段4图像任务样本浏览与提示词核查界面}
    \label{fig:app-stage4-data-browser}
\end{figure}

\autoref{fig:app-stage4-data-browser} 展示了阶段4的图像任务浏览界面。左侧为任务清单和任务组合统计，右侧同步展示样本摘要、系统提示词、用户提示词、候选选项和目标框字段。该界面用于核查参考图像、查询图像、候选答案、提示词协议和目标框监督是否一致，尤其适合检查自我状态认知任务和环境情况认知任务是否遵循同一输出契约。

\begin{figure}[!htbp]
    \centering
    \includegraphics[width=\textwidth]{fig_app_stage4_experiments.png}
    \caption{阶段4图像任务实验调度工作台界面}
    \label{fig:app-stage4-experiments}
\end{figure}

\autoref{fig:app-stage4-experiments} 展示了阶段4图像任务的实验配置界面。该界面记录样本清单、上传分辨率、JPEG 压缩质量、并发数、超时限制和待评测模型等实验条件，并在同一视图中呈现实验队列状态。

\begin{figure}[!htbp]
    \centering
    \includegraphics[width=\textwidth]{fig_app_stage4_metrics.png}
    \caption{阶段4图像任务指标汇总工作台界面}
    \label{fig:app-stage4-metrics}
\end{figure}

\autoref{fig:app-stage4-metrics} 展示了阶段4图像任务指标汇总界面。该界面从样本级输出中聚合解析成功率、选项准确率、BBox Acc@0.5、BBox mIoU 和平均延迟，并进一步按照认知分支、任务类型和难度分组。该界面对应正文中图像任务实验表格的指标来源，也便于发现某一任务分支或难度设置下的系统性弱点。

\subsection{四阶段构建流程伪代码}

\begin{verbatim}
Input: AerialVLN scenes and task configuration
Output: UAV-DualCog benchmark with image/video evidence labels

for each selected scene do
    Stage 1: scan scene and collect RGB / semantics / LiDAR / poses
    fuse observations into a scene-level semantic point cloud

    Stage 2: cluster landmark candidates from semantic point cloud
    render multi-view landmark candidates
    conduct human review and semantic labeling
    export validated landmark assets

    Stage 3: bind landmarks with flight behavior templates
    generate trajectories and first-person videos
    record temporal evidence for actions and landmark visibility
    export self-aware and environment-aware video tasks

    Stage 4: sample landmark-centered reference view
    sample UAV observation view or future-view candidates
    generate image QA tasks and normalized bbox evidence
end for
\end{verbatim}

阶段3的关键在于为同一条轨迹同时建立高层行为语义与底层时间证据。阶段4的关键在于同步生成选项答案和目标框证据，而不是先生成问题再事后补充定位框。因此，UAV-DualCog 的任务构建天然具备“语义答案 + 显式证据”的联合评测属性。
\section{评测协议补充材料}
\label{app:evaluation-protocol}
\label{app:prompt-protocol}

UAV-DualCog 的评测协议由系统提示词、用户提示词模板、结构化输出协议和指标解析规则共同组成。系统提示词统一限定模型扮演无人机智能体、明确输入媒介和禁止额外解释；用户提示词注入地标描述、候选选项、参考图像、查询图像或视频语境；输出协议要求模型返回严格 JSON，便于解析器稳定计算语义指标和证据指标。

\begin{table}[htbp]
    \centering
    \caption{评测输出字段与指标映射}
    \label{tab:app-output-metric-map}
    \footnotesize
    \begin{tabularx}{\textwidth}{@{}U{3.0cm}U{3.4cm}X@{}}
        \toprule
        输出字段 & 适用任务 & 驱动指标 \\
        \midrule
        \makecell[c]{\texttt{answer\_}\\\texttt{option\_id}} & 图像任务 & Option Accuracy \\
        \texttt{bbox\_xyxy\_norm} & 图像任务 & BBox Acc@0.5、BBox mIoU \\
        \makecell[c]{\texttt{answers[]}\\\texttt{.intervals\_sec}} & 自我状态认知视频任务 & 行为识别主指标、Temporal F1@0.5、Mean tIoU \\
        \texttt{visible\_count} & 环境情况认知视频任务 & Count Accuracy、Count Within1 \\
        \makecell[c]{\texttt{visible\_}\\\texttt{intervals\_sec}} & 环境情况认知视频任务 & Segment F1@0.5、Mean tIoU \\
        \bottomrule
    \end{tabularx}
\end{table}

\begin{table}[htbp]
    \centering
    \caption{六类任务提示词协议摘要}
    \label{tab:app-prompt-protocol}
    \footnotesize
    \setlength{\tabcolsep}{4pt}
    \begin{tabularx}{\textwidth}{@{}U{2.8cm}U{2.5cm}X U{2.8cm}@{}}
        \toprule
        任务 & 输入媒介 & 核心提示信息 & 必要输出字段 \\
        \midrule
        地标参照位置推理 & 参考图像 + 查询图像 & 地标中心参考视图、当前无人机观测图、相对位置选项 & \makecell[c]{\texttt{answer\_option\_id}\\\texttt{bbox\_xyxy\_norm}} \\
        未来观察预测 & 参考图像 + 候选图像 & 当前观察、轨道运动描述、多张候选未来观察图 & \makecell[c]{\texttt{answer\_option\_id}\\\texttt{bbox\_xyxy\_norm}} \\
        自我参照位置推理 & 查询图像 & 当前第一视角观察、地标描述、方位选项 & \makecell[c]{\texttt{answer\_option\_id}\\\texttt{bbox\_xyxy\_norm}} \\
        地标驱动动作决策 & 查询图像 & 当前第一视角观察、目标地标描述、动作选项 & \makecell[c]{\texttt{answer\_option\_id}\\\texttt{bbox\_xyxy\_norm}} \\
        飞行行为识别与时间定位 & 视频 & 第一视角飞行视频、复合或原子行为候选选项 & \makecell[c]{\texttt{answers}\\\texttt{intervals\_sec}} \\
        地标可见次数与时间定位 & 视频 + 参考图像 & 目标地标参考图、第一视角飞行视频、可见性问题 & \makecell[c]{\texttt{visible\_count}\\\texttt{visible\_}\\\texttt{intervals\_sec}} \\
        \bottomrule
    \end{tabularx}
\end{table}

\subsection{图像任务完整中文提示词模板}

下列模板为实际英文评测提示词的中文等价版本。模板中的占位字段由样本构建程序根据具体任务自动替换，主要包括：
\begin{itemize}
    \item \texttt{reference\_object\_view}：参考立面或参考视角字段；
    \item \texttt{landmark\_description}：目标地标文本描述字段；
    \item \texttt{behavior\_instance}：飞行动作或轨迹变化字段；
    \item \texttt{options\_text}：候选选项字段。
\end{itemize}

\paragraph{地标参照位置推理。}
\textbf{系统提示词：} 你是一个 UAV agent。你将获得一张包含目标地标的参考视图和一张当前第一视角图像。图像1为参考视图，显示目标地标及其目标框；图像2为当前视图。图像1中可见的立面定义了地标中心坐标系下的参考方向。你必须根据该参考方向判断自己在图像2中相对于地标的位置，并同时定位图像2中的目标地标。只允许返回一个合法 JSON 对象，不得输出 Markdown、解释或额外文本；只能使用用户提示词中给出的选项编号；目标框必须使用归一化 \([x_1,y_1,x_2,y_2]\) 坐标。
\begin{verbatim}
{"answer_option_id":"A","bbox_xyxy_norm":[0.1,0.2,0.3,0.4]}
\end{verbatim}
\textbf{用户提示词：} 图像1显示地标中心坐标系中 \texttt{\{reference\_object\_view\}} 立面的 \texttt{\{landmark\_description\}}。基于这一参考，图像2中你相对于该地标处于什么位置？请选择一个选项，并返回图像2中该地标的归一化目标框。选项为：\texttt{\{options\_text\}}。解析字段为 \texttt{answer\_option\_id} 和 \texttt{bbox\_xyxy\_norm}。

\paragraph{未来观察预测。}
\textbf{系统提示词：} 你是一个 UAV agent。你将获得一张包含目标地标的参考视图和四张候选未来视图。图像1为参考视图，显示目标地标及其目标框；图像2到图像5为围绕同一地标的候选未来视图。图像1中可见的立面定义了地标中心坐标系下的参考方向。你必须根据给定轨道运动选择唯一正确的未来视图，并在所选候选图像中定位目标地标。只允许返回一个合法 JSON 对象，不得输出 Markdown、解释或额外文本；目标框必须使用归一化坐标。
\begin{verbatim}
{"answer_option_id":"A","bbox_xyxy_norm":[0.1,0.2,0.3,0.4]}
\end{verbatim}
\textbf{用户提示词：} 图像1显示地标中心坐标系中 \texttt{\{reference\_object\_view\}} 立面的 \texttt{\{landmark\_description\}}。执行 \texttt{\{behavior\_instance\}} 后，哪一张候选图像是对应的未来视图？请选择一个选项，并返回所选图像中该地标的归一化目标框。选项为：A. 图像2；B. 图像3；C. 图像4；D. 图像5。解析字段为 \texttt{answer\_option\_id} 和 \texttt{bbox\_xyxy\_norm}。

\paragraph{自我参照位置推理。}
\textbf{系统提示词：} 你是一个 UAV agent。你将获得目标地标参考视图和当前第一视角图像。图像1为参考视图，图像2为当前视图。你需要判断目标地标在图像2中相对于自己前进方向的位置，并同时定位该地标。只允许返回一个合法 JSON 对象，不得输出 Markdown、解释或额外文本；只能使用用户提示词中给出的选项编号；目标框必须使用归一化坐标。
\begin{verbatim}
{"answer_option_id":"A","bbox_xyxy_norm":[0.1,0.2,0.3,0.4]}
\end{verbatim}
\textbf{用户提示词：} 图像1显示地标中心坐标系中 \texttt{\{reference\_object\_view\}} 立面的 \texttt{\{landmark\_description\}}。在图像2中，该地标相对于你的前进方向位于哪里？请选择一个选项，并返回图像2中该地标的归一化目标框。选项为：\texttt{\{options\_text\}}。解析字段为 \texttt{answer\_option\_id} 和 \texttt{bbox\_xyxy\_norm}。

\paragraph{地标驱动动作决策。}
\textbf{系统提示词：} 你是一个 UAV agent。你将获得目标地标参考视图和当前第一视角图像。你需要判断为了接近该地标应向哪个方向飞行，并同时定位当前图像中的目标地标。只允许返回一个合法 JSON 对象，不得输出 Markdown、解释或额外文本；只能使用用户提示词中给出的选项编号；目标框必须使用归一化坐标。
\begin{verbatim}
{"answer_option_id":"A","bbox_xyxy_norm":[0.1,0.2,0.3,0.4]}
\end{verbatim}
\textbf{用户提示词：} 图像1显示地标中心坐标系中 \texttt{\{reference\_object\_view\}} 立面的 \texttt{\{landmark\_description\}}。在图像2中，为了接近该地标，你应该向哪个方向飞行？请选择一个选项，并返回图像2中该地标的归一化目标框。选项为：\texttt{\{options\_text\}}。解析字段为 \texttt{answer\_option\_id} 和 \texttt{bbox\_xyxy\_norm}。

\subsection{视频任务完整中文提示词模板}

\paragraph{复合飞行行为识别。}
\textbf{系统提示词：} 你是一个 UAV agent。请观看自己的第一视角飞行视频，识别视频中曾经执行过的高层飞行动作。用户提示词会提供候选选项。只允许返回一个合法 JSON 对象，不得输出 Markdown、解释或额外文本；只能使用用户提示词中出现的 \texttt{option\_id}；每个答案项必须包含选项编号及其对应秒级时间区间；若多个候选动作确实出现，需要全部返回。
\begin{verbatim}
{"answers":[{"option_id":"A","intervals_sec":[[0.0,1.2]]}]}
\end{verbatim}
\textbf{用户提示词：} 在该第一视角飞行视频中，你曾经执行过以下哪些动作？请选择一个或多个选项，并返回对应的秒级时间区间。选项为：\texttt{\{options\_text\}}。解析字段为 \texttt{answers[].option\_id} 和 \texttt{answers[].intervals\_sec}。

\paragraph{原子飞行行为识别。}
\textbf{系统提示词：} 你是一个 UAV agent。请观看自己的第一视角飞行视频，识别视频中执行过的低层原子飞行动作。用户提示词会提供候选选项。只允许返回一个合法 JSON 对象，不得输出 Markdown、解释或额外文本；只能使用用户提示词中出现的 \texttt{option\_id}；每个答案项必须包含选项编号及其对应秒级时间区间；若多个候选动作确实出现，需要全部返回。
\begin{verbatim}
{"answers":[{"option_id":"A","intervals_sec":[[2.1,4.6]]}]}
\end{verbatim}
\textbf{用户提示词：} 在该第一视角飞行视频中，你执行了以下哪一个动作？请选择一个选项，并返回对应的秒级时间区间；若多个候选动作确实出现，则需要全部返回。选项为：\texttt{\{options\_text\}}。解析字段为 \texttt{answers[].option\_id} 和 \texttt{answers[].intervals\_sec}。

\paragraph{地标可见次数与时间定位。}
\textbf{系统提示词：} 你是一个 UAV agent。请观看自己的第一视角飞行视频，并结合一张目标地标参考图判断该地标何时可见。你将获得一个第一视角飞行视频、一张显示目标地标及目标框的参考图像，以及该地标的简短文本描述。只允许返回一个合法 JSON 对象，不得输出 Markdown、解释或额外文本；需要统计目标地标分几次可见，并返回完整可见时间区间。
\begin{verbatim}
{"visible_count":2,"visible_intervals_sec":[[0.0,9.1],[17.7,26.8]]}
\end{verbatim}
\textbf{用户提示词：} 参考图显示 \texttt{\{landmark\_description\}}。该地标在你的第一视角飞行视频中出现了几次？请返回完整可见时间区间，单位为秒。解析字段为 \texttt{visible\_count} 和 \texttt{visible\_intervals\_sec}。

\subsection{解析与有效性规则}

所有任务均要求严格 JSON 输出。包含 Markdown 代码块、额外自然语言解释、字段缺失或字段类型错误的结果会被视为格式无效。图像任务只有在选项与目标框均可解析时才进入完整评分；视频任务则将语义识别或可见次数与时间区间质量分开计算，从而避免把“识别正确但证据错误”的输出误判为完整成功。
\section{训练集及构建补充材料}
\label{app:train-dataset-construction}

UAV-DualCog-Train-V1 面向第4章能力优化实验构建，仅使用与测试集场景隔离的训练场景。其目标不是扩大测试集，而是将第3章暴露出的“图像语义判断与空间证据脱节”问题转化为可监督、可奖励的训练记录。

\begin{table}[htbp]
    \centering
    \caption{UAV-DualCog-Train-V1 训练数据字段设计}
    \label{tab:app-train-fields}
    \small
    \begin{tabularx}{\textwidth}{@{}U{3.1cm}U{3.0cm}X@{}}
        \toprule
        字段 & 类型 & 作用 \\
        \midrule
        \texttt{prompt} & 对话文本 & 保存与评测协议一致的系统提示词和用户提示词 \\
        \texttt{completion} & JSON 字符串 & 保存监督目标，包括 \texttt{answer\_option\_id} 和 \texttt{bbox\_xyxy\_norm} \\
        \texttt{images} & 图像路径列表 & 按评测输入顺序保存真实图像路径 \\
        \texttt{task\_type} & 字符串 & 区分四类图像训练任务 \\
        \texttt{difficulty} & 字符串 & 区分两档问题难度 \\
        \makecell[c]{\texttt{answer\_bbox}\\\texttt{\_xyxy\_1000}} & 数值列表 & 供奖励函数复用的目标框辅助字段 \\
        \bottomrule
    \end{tabularx}
\end{table}

\begin{table}[H]
    \centering
    \caption{UAV-DualCog-Train-V1 与测试集场景隔离划分}
    \label{tab:app-scene-split}
    \footnotesize
    \setlength{\tabcolsep}{3pt}
    \begin{tabularx}{\textwidth}{@{}U{1.8cm}V{4.5cm}U{2.0cm}U{1.5cm}X@{}}
        \toprule
        数据部分 & 场景 ID & \makecell[c]{场景面积\\（m\textsuperscript{2}）} & 地标数量 & 用途说明 \\
        \midrule
        训练集场景 & \makecell[l]{\texttt{env\_1}, \texttt{env\_3}, \texttt{env\_5}\\\texttt{env\_15}, \texttt{env\_18}, \texttt{env\_22}} & 1,596,511 & 234 & 用于构建 UAV-DualCog-Train-V1 图像训练样本，参与 SFT 与联合 GRPO 训练。 \\
        测试集场景 & \makecell[l]{\texttt{env\_7}, \texttt{env\_8}, \texttt{env\_9}\\\texttt{env\_10}, \texttt{env\_11}, \texttt{env\_13}\\\texttt{env\_16}, \texttt{env\_17}, \texttt{env\_20}\\\texttt{env\_21}, \texttt{env\_23}, \texttt{env\_24}} & 3,443,721 & 512 & 用于 UAV-DualCog-V1 图像与视频任务评测，不参与训练数据构建。 \\
        \bottomrule
    \end{tabularx}
\end{table}

其中，测试集场景面积与地标数量来自 UAV-DualCog-V1 发布统计，训练集场景面积由完整资产池面积扣除测试场景面积得到，训练集地标数量来自 UAV-DualCog-Train-V1 构建统计。

\begin{table}[htbp]
    \centering
    \caption{UAV-DualCog-Train-V1 数据分布与完整性补充统计}
    \label{tab:app-train-stat}
    \footnotesize
    \setlength{\tabcolsep}{4pt}
    \begin{tabularx}{\textwidth}{@{}U{4.1cm}U{2.2cm}U{4.1cm}U{2.2cm}@{}}
        \toprule
        统计项 & 数值 & 统计项 & 数值 \\
        \midrule
        有效训练地标数 & 234 & 全量样本数 & 11232 \\
        训练集样本数 & 10104 & 验证集样本数 & 1128 \\
        四类任务样本数 & 每类 2808 & 双认知分组样本数 & 各 5616 \\
        两档难度样本数 & 各 5616 & 每个任务-难度组合 & 1404 \\
        非法目标框数量 & 0 & 图像路径存在率 & 100\% \\
        目标框面积均值 & 0.0645 & 训练/验证比例 & 89.96\% : 10.04\% \\
        \bottomrule
    \end{tabularx}
\end{table}

\subsection{训练数据准备伪代码}

\begin{verbatim}
Input: stage4 manifests from six training scenes
Output: all / train / valid JSONL and integrity report

for each scene manifest do
    load image-task samples
    build benchmark-aligned system/user prompt
    collect real image paths in benchmark input order
    build assistant completion as JSON:
        {"answer_option_id": "...",
         "bbox_xyxy_norm": [...]}
    cache normalized bbox and bbox_1000 for reward computation
end for

bucket samples by (task_type, difficulty)
sort each bucket by sample id
assign every 10th sample to valid split
assign remaining samples to train split
dump JSONL files and integrity report
\end{verbatim}

该构建方式保证训练样本在任务类型、难度等级和双认知分组上保持均衡，同时保证训练集与测试集在场景层面隔离。
\section{能力优化实现补充材料}
\label{app:optimization-entry}

本附录补充第4章能力优化训练环境、运行记录、实现入口、训练流程伪代码、完整性核查逻辑和奖励函数关键片段。

\subsection{训练环境与运行记录}

\begin{table}[htbp]
    \centering
    \caption{能力优化实验训练环境}
    \label{tab:app-train-hardware}
    \small
    \begin{tabularx}{\textwidth}{@{}U{3.0cm}>{\centering\arraybackslash}X@{}}
        \toprule
        资源项 & 配置 \\
        \midrule
        CPU & \(2\times\) AMD EPYC 7452，合计 128 vCPUs \\
        内存 & 503 GiB \\
        GPU & \(4\times\) NVIDIA GeForce RTX 4090，单卡 24 GB 显存 \\
        训练框架 & PyTorch、Transformers、PEFT、TRL \\
        训练方式 & 单机多卡多模态 LoRA-SFT + 联合 GRPO \\
        \bottomrule
    \end{tabularx}
\end{table}

\begin{table}[htbp]
    \centering
    \caption{正式训练完成情况与核心过程指标}
    \label{tab:app-train-run-summary}
    \footnotesize
    \setlength{\tabcolsep}{4pt}
    \begin{tabularx}{\textwidth}{@{}U{1.9cm}U{1.6cm}U{1.6cm}U{2.0cm}X@{}}
        \toprule
        阶段 & 完成步数 & 运行时间/h & 过程指标 & 说明 \\
        \midrule
        SFT & 237/237 & 4.05 & \makecell[c]{train loss\\0.6384} & 完成 3 个 epoch，验证集 \texttt{eval\_loss}=0.5666，图像长边上限为 720。 \\
        联合 GRPO & 600/600 & 10.72 & \makecell[c]{tail-10 reward\\2.4016} & 奖励权重为 \((1.0,1.0,1.0)\)，后段奖励保持稳定，导出可评测权重。 \\
        \bottomrule
    \end{tabularx}
\end{table}

\begin{table}[htbp]
    \centering
    \caption{能力优化主要入口与作用说明}
    \label{tab:app-optimization-entry}
    \scriptsize
    \begin{tabularx}{\textwidth}{@{}U{2.6cm}V{4.1cm}X@{}}
        \toprule
        脚本入口 & 主要输入输出 & 作用说明 \\
        \midrule
        \makecell[c]{\texttt{prepare\_train}\\\texttt{\_dataset.py}} & 输入 Stage 4 manifest；输出 train/valid/all JSONL 与完整性报告 & 将 UAV-DualCog-Train-V1 图像任务转换为统一多模态训练记录。 \\
        \makecell[c]{\texttt{run\_sft}\\\texttt{\_lora.py}} & 输入 train/valid JSONL 与 Qwen 3.5 4B；输出 LoRA adapter、训练指标与合并模型 & 实现多模态监督微调，使真实图像参与训练并学习结构化 JSON 输出。 \\
        \makecell[c]{\texttt{run\_grpo}\\\texttt{\_light.py}} & 输入训练 JSONL 与 SFT 检查点；输出 GRPO adapter、奖励指标与合并模型 & 实现联合 GRPO，在同一奖励函数中优化格式、答案和目标框质量。 \\
        \makecell[c]{\texttt{run\_full}\\\texttt{\_pipeline.sh}} & 管理数据快照、日志、检查点和两阶段调度 & 串联数据准备、SFT 与联合 GRPO，形成单机多卡可复现训练流程。 \\
        \makecell[c]{\texttt{check\_training}\\\texttt{\_run.py}} & 输入运行目录；输出完整性报告 & 核查计划文件、日志、checkpoint、adapter 和合并模型产物是否完整。 \\
        \makecell[c]{\texttt{evaluate}\\\texttt{\_transfer.py}} & 输入原始模型/优化模型权重与 benchmark 场景；输出评测结果与增益表 & 触发 Stage 3/4 评测并生成优化模型对比结果。 \\
        \bottomrule
    \end{tabularx}
\end{table}

\subsection{两阶段训练流程伪代码}

\begin{verbatim}
Input: UAV-DualCog-Train-V1 JSONL, Qwen 3.5 4B base model
Output: SFT checkpoint and GRPO checkpoint

Stage 1: SFT
    load train / valid JSONL
    load real RGB images on demand
    encode text and images with AutoProcessor
    train LoRA adapters with completion-only loss
    export adapter and merged model

Stage 2: Joint GRPO
    initialize from SFT checkpoint
    load real RGB images on demand
    generate JSON completions under multimodal model
    compute format reward + answer reward + bbox reward
    update LoRA adapters with group-relative optimization
    export adapter and merged model
\end{verbatim}

\subsection{正式训练完整性核查逻辑}

正式训练结束后，系统对运行目录进行完整性核查，重点确认数据快照、训练计划、日志文件、指标文件、checkpoint、LoRA adapter、合并模型和评测入口是否完整。

\begin{verbatim}
Input: training run directory
Output: integrity JSON / Markdown report

check run root exists
for each stage in [01_sft_full, 02_grpo_joint] do
    verify plan file, metrics file, log file exist
    verify adapter and merged_model exist
    count checkpoints and locate latest step
    parse trainer_state tail summary
end for

if both stages complete and merged models exist then
    mark overall_complete = true
\end{verbatim}

\subsection{奖励函数关键代码片段}
\label{app:reward-code}

下列代码片段节选自奖励实现，展示本文如何优先解析 JSON 结构化输出，并将格式合法性、答案正确性和目标框一致性组合为统一奖励。
\begin{verbatim}
def reward_format(text: str) -> float:
    return 1.0 if (
        parse_answer(text) is not None
        and parse_bbox_norm(text) is not None
    ) else 0.0

def reward_mcq(pred_text: str, gt_answer: str) -> float:
    pred = parse_answer(pred_text)
    if pred is None:
        return 0.0
    return 1.0 if pred == str(gt_answer).strip().upper() else 0.0

def reward_bbox(pred_text: str, gt_bbox_norm):
    pred_bbox = parse_bbox_norm(pred_text)
    if pred_bbox is None:
        return 0.0
    giou = giou_norm(pred_bbox, gt_bbox_norm)
    return max(0.0, min(1.0, (giou + 1.0) * 0.5))

def total_reward(pred_text, gt_answer, gt_bbox_norm, weights):
    return (
        weights.w_format * reward_format(pred_text)
        + weights.w_mcq * reward_mcq(pred_text, gt_answer)
        + weights.w_bbox * reward_bbox(pred_text, gt_bbox_norm)
    )
\end{verbatim}

\end{thesisappendix}
